% !TEX root = ../main.tex
\begin{figure}[H]
  \centering
  \resizebox{\textwidth}{!}{%
    \begin{tikzpicture}[
        font=\footnotesize,
        bloque/.style={draw=diagNodoBorde, fill=diagNodo, align=left,
            inner sep=7pt, rounded corners=2pt},
        titulo/.style={font=\footnotesize\bfseries, anchor=north west,
            text=diagGrupoBorde!55!black}
      ]
      \node[bloque, text width=9.2cm] (enc) {%
        \textbf{Encabezado autocontenido} \hfill \textit{generado, no extraído}\\[3pt]
        <<Automática industrial>>, asignatura obligatoria de
        6 ECTS impartida en 8 titulaciones: Grado en Ingeniería Electrónica
        Industrial; Grado en Ingeniería Eléctrica; Grado en Ingeniería Mecánica;
        Grado en Ingeniería de Organización Industrial; y los dobles grados en
        Electrónica Industrial y Mecánica, Eléctrica y Electrónica Industrial,
        Eléctrica y Mecánica, y Mecánica y Organización Industrial.};
      \node[bloque, text width=9.2cm, below=4mm of enc.south west,
        anchor=north west] (cue) {%
        \textbf{Cuerpo} \hfill \textit{texto de la guía docente}\\[3pt]
        Teoría. Tema 1: Introducción. 1.1 ¿Qué es la automática (industrial)?
        1.2 Sistemas (en automática).\dots};

      \node[bloque, text width=5.6cm, right=9mm of enc.north east,
        anchor=north west] (met) {%
        \textbf{Metadatos}\\[3pt]
        \texttt{origen}: guia\\
        \texttt{nombre}: Automática industrial\\
        \texttt{grados}: [8 titulaciones]\\
        \texttt{codigos}: [13112002, 13512002,\\ \hphantom{\texttt{codigos}: [}13412001, 13012001,\\ \hphantom{\texttt{codigos}: [}13912002, 13712003,\\ \hphantom{\texttt{codigos}: [}13612002, 13812001]\\
        \texttt{chunk\_index}: 1\\
        \texttt{total\_chunks}: 14};

      \begin{scope}[on background layer]
        \node[grupoDiag, rounded corners=3pt, fit=(enc)(cue)] (gtxt) {};
        \node[grupoDiag, rounded corners=3pt, fit=(met)] (gmet) {};
      \end{scope}
      \node[titulo, anchor=south west] (ltxt)
      at ([yshift=1.4mm]gtxt.north west) {campo \texttt{texto}: lo único que se convierte en vector};
      \node[titulo, anchor=south west] (lmet)
      at ([yshift=1.4mm]gmet.north west) {acompañan al vector, no se buscan};

      \begin{scope}[on background layer]
        \node[draw=diagNodoBorde, thick, rounded corners=4pt, inner sep=9pt,
          fit=(gtxt)(gmet)(ltxt)(lmet)] (marco) {};
      \end{scope}
      \node[font=\footnotesize\bfseries, anchor=south west]
      at ([yshift=1.8mm]marco.north west)
      {Segundo fragmento de 14 de <<Automática industrial>>};
    \end{tikzpicture}%
  }
  \caption[Estructura de un fragmento (\textit{chunk})]{Estructura de un
    fragmento de la colección, sobre un caso real de \texttt{chunks.json}. El
    encabezado se antepone al cuerpo para que el fragmento siga teniendo sentido
    al recuperarse aislado; los campos \texttt{grados} y \texttt{codigos} son
    listas porque esta guía está deduplicada y se imparte en
    ocho titulaciones. Fuente: Elaboración propia.}
  \label{fig:estructura_chunk}
\end{figure}
